\documentclass[12pt]{article}
\usepackage[utf8]{inputenc}
\usepackage[T1]{fontenc}
\usepackage{hyperref}
\title{Suba Language Resources for K-12 Curriculum Development}
\author{Compiled by Claude Sonnet 5.5 (claude-sonnet-5-5) for Andr\'as Kornai}
\date{}
\begin{document}
\maketitle

\section{Introduction}
Suba (Olusuba, Kisuba; ISO 639-3 sxb) is a Bantu language of Lake Victoria in Kenya. It is spoken on the eastern shore and on the islands of Mfangano and Rusinga in Homa Bay County. It is closely related to Ganda and has borrowed heavily from Dholuo. Reports say that many children on the islands do not acquire it in infancy, so revitalization is a live issue.

\section{Primary Educational Resources}
Search results state that Suba was introduced as a subject in Kenyan primary schools in the 1990s. No specific textbooks or primers were found.

\section{Linguistic and Reference Materials}
\begin{itemize}
\item Wikipedia entry on the Suba language: \url{https://en.wikipedia.org/wiki/Suba_language}
\item Omniglot page on Suba (Olusuba): \url{https://omniglot.com/writing/suba.htm}
\end{itemize}
No dictionary was found in the searches.

\section{Religious and Community Resources}
\begin{itemize}
\item Wikipedia entry on the Suba people of Kenya, which mentions the Suba Language Project and the Rusinga Cultural Festival: \url{https://en.wikipedia.org/wiki/Suba_people_(Kenya)}
\item Nation article on the push to save the Suba language: \url{https://www.nation.africa/kenya/counties/homa-bay/-push-to-save-dying-suba-language-now-gathers-momentum-4128100}
\item Search results report a Suba New Testament published in 2010. I did not locate it online.
\end{itemize}

\section{Digital and Audio Resources}
\begin{itemize}
\item Global Recordings Network Suba page: \url{https://www.globalrecordings.net/es/suba-distribution-2026}
\item Endangered Languages Project entry: \url{https://endangeredlanguages.com/lang/4789?hl=fr}
\end{itemize}

\section{Curriculum Development Context}
Suba is reported as taught as a subject in some schools, but it is not used as a medium of instruction. A curriculum would build on the existing subject syllabus, the Suba Language Project and community organizations, and would need new teaching materials and a dictionary. Existing materials appear limited.
\end{document}